\problem{4.63}{Normalization of all spherical harmonics}
Write $Y_\ell^m=K_{\ell m}
P_\ell^m(\cos\theta)\ee^{\ii m\phi}$.
The differential ladder operator and
the identity for associated Legendre
functions imply
\begin{equation}
L_+Y_\ell^m=\sh K_{\ell m}
P_\ell^{m+1}(\cos\theta)
\ee^{\ii(m+1)\phi}.
\end{equation}
But the normalized angular-momentum rule is
$L_+Y_\ell^m=
\sh\sqrt{(\ell-m)(\ell+m+1)}Y_\ell^{m+1}$.
Consequently
\begin{equation}
\boxed{\frac{K_{\ell,m+1}}{K_{\ell m}}
=\frac1{\sqrt{(\ell-m)(\ell+m+1)}}}.
\end{equation}
Iterating from $m=0$ gives
\begin{equation}
K_{\ell m}=K_{\ell0}
\sqrt{\frac{(\ell-m)!}{(\ell+m)!}}.
\end{equation}
Orthogonality of ordinary Legendre polynomials
requires
$1=2\pi K_{\ell0}^2\int_{-1}^1P_\ell^2\dd x
=4\pi K_{\ell0}^2/(2\ell+1)$.
Thus
\begin{equation}
\boxed{K_{\ell m}=
\sqrt{\frac{2\ell+1}{4\pi}
\frac{(\ell-m)!}{(\ell+m)!}}}.
\end{equation}
The phase is fixed by the Condon--Shortley
convention built into $P_\ell^m$.
Negative $m$ follows from
$Y_\ell^{-m}=(-1)^mY_\ell^{m*}$.

\problem{4.64}{Combined orbital and spin hydrogen measurements}
Consider the normalized state
\begin{equation}
\ket\Psi=\sqrt{\frac13}\ket{\ell=1,m_\ell=0}
\ket{\uparrow}+
\sqrt{\frac23}\ket{1,1}\ket{\downarrow},
\end{equation}
with the common normalized $R_{21}(r)$.
Both components have $\ell=1$ and $s=1/2$.
Therefore
\begin{equation}
\boxed{L^2=2\sh^2\;(P=1),\qquad
S^2=\tfrac34\sh^2\;(P=1)}.
\end{equation}
For $L_z$ the outcomes are $0$ with $P=1/3$
and $\sh$ with $P=2/3$; for $S_z$, the outcomes
$\sh/2,-\sh/2$ have $P=1/3,2/3$.
Both components have $m_j=m_\ell+m_s=1/2$,
so $\boxed{J_z=\sh/2\;(P=1)}$.
For $j=3/2,m_j=1/2$ the CG vector is
$\sqrt{2/3}\ket{0,\uparrow}+
\sqrt{1/3}\ket{1,\downarrow}$.
Projecting gives amplitude
$\sqrt2/3+\sqrt2/3=2\sqrt2/3$, hence
\begin{equation}
\boxed{P(J^2=\tfrac{15}{4}\sh^2)=\tfrac89,\qquad
P(J^2=\tfrac34\sh^2)=\tfrac19}.
\end{equation}
For position measurement the spin alternatives
are orthogonal, so no spatial interference term:
\begin{align}
|\Psi(\mathbf r)|^2&=|R_{21}(r)|^2
\left[\tfrac13|Y_1^0|^2+
\tfrac23|Y_1^1|^2\right]\\
&=\boxed{\frac{|R_{21}(r)|^2}{4\pi}},
\end{align}
which is isotropic. The probability per unit
radius for spin up is
\begin{equation}
\boxed{P_+(r)=\frac13r^2|R_{21}(r)|^2}.
\end{equation}

\problem{4.65}{Three spin halves: one quadruplet and two doublets}
Use the shorthand $u=\uparrow$, $d=\downarrow$.
The totally symmetric spin-$3/2$ quadruplet is
\begin{align}
\ket{\tfrac32,\tfrac32}&=\ket{uuu},\\
\ket{\tfrac32,\tfrac12}&=
(\ket{uud}+\ket{udu}+\ket{duu})/\sqrt3,\\
\ket{\tfrac32,-\tfrac12}&=
(\ket{udd}+\ket{dud}+\ket{ddu})/\sqrt3,\\
\ket{\tfrac32,-\tfrac32}&=\ket{ddd}.
\end{align}
The first spin-$1/2$ doublet is obtained by
placing particles 1 and 2 in a singlet:
\begin{align}
\ket{\tfrac12,\tfrac12}_{1}
&=(\ket{udu}-\ket{duu})/\sqrt2,\\
\ket{\tfrac12,-\tfrac12}_{1}
&=(\ket{udd}-\ket{dud})/\sqrt2.
\end{align}
A second independent doublet, orthogonal
to both the first doublet and the quadruplet,
is
\begin{align}
\ket{\tfrac12,\tfrac12}_{2}
&=(2\ket{uud}-\ket{udu}-\ket{duu})/\sqrt6,\\
\ket{\tfrac12,-\tfrac12}_{2}
&=(\ket{udd}+\ket{dud}-2\ket{ddu})/\sqrt6.
\end{align}
Each state has unit norm, the two doublets
are mutually orthogonal, and total
ladder operators connect the $\pm1/2$
members within each doublet. Dimension
count gives $4+2+2=8=2^3$.

\problem{4.66}{Saturation of the spin uncertainty bound}
Let $\chi=(a,b)^T$, $\|\chi\|=1$.
Remove an overall phase to make $a\ge0$ real.
Define the Bloch components
$n_x=2a\operatorname{Re}b$,
$n_y=2a\operatorname{Im}b$,
$n_z=|a|^2-|b|^2$.
A pure spinor obeys $n_x^2+n_y^2+n_z^2=1$.
Since $\sigma_{S_i}^2=
(\sh^2/4)(1-n_i^2)$,
\begin{align}
\sigma_{S_x}^2\sigma_{S_y}^2
-\frac{\sh^2}{4}\avg{S_z}^2
&=\frac{\sh^4}{16}
[(1-n_x^2)(1-n_y^2)-n_z^2]\\
&=\frac{\sh^4}{16}n_x^2n_y^2.
\end{align}
Equality therefore holds if and only if
$n_xn_y=0$. With $a$ real, this means
\begin{equation}
\boxed{\operatorname{Re}b=0
\quad\text{or}\quad\operatorname{Im}b=0},
\end{equation}
so $b$ can be chosen purely imaginary
or purely real. The special case $a=0$
automatically saturates the inequality.

\problem{4.67}{Antiferromagnetic frustration on triangles and squares}
For three spins $\mathbf S_i$ of magnitude
$s_i=1/2$, define $\mathbf S=\sum_i\mathbf S_i$.
Then
\begin{equation}
S^2=\sum_iS_i^2+
2(S_1\cdot S_2+S_2\cdot S_3+S_3\cdot S_1).
\end{equation}
Hence for the triangular Hamiltonian
\begin{equation}
\boxed{H_\triangle=
\frac J2(S^2-\tfrac94\sh^2)}.
\end{equation}
Possible total spins are $3/2$ and $1/2$;
the latter occurs in two independent doublets.
Their energies are
$E_{1/2}=-3J\sh^2/4$ and
$E_{3/2}=3J\sh^2/4$.
Thus
\begin{equation}
\boxed{E_{\rm ground}^{\triangle}=-\tfrac34J\sh^2,
\qquad d_{\rm ground}=4}.
\end{equation}
For a square, group opposite vertices:
$\mathbf A=\mathbf S_1+\mathbf S_3$,
$\mathbf B=\mathbf S_2+\mathbf S_4$.
Then $H_\square=J\mathbf A\cdot\mathbf B$
and
\begin{equation}
\boxed{H_\square=
\frac J2(S^2-A^2-B^2)}.
\end{equation}
Each opposite-vertex pair can have
spin $0$ or $1$. Choosing $A=B=1$
and combining to total $S=0$ minimizes
the energy:
\begin{equation}
\boxed{E_{\rm ground}^{\square}=-2J\sh^2},
\end{equation}
with a unique singlet ground state.

\problem{4.68}{Numerical radial hydrogen in a large spherical wall}
For a Coulomb center with a hard wall
$u(b)=0$, discretize
$u(r)=rR(r)$ at $r_j=j\Delta r$,
$\Delta r=b/(N+1)$, and put
$\beta=\Delta r/\aB=b/[(N+1)\aB]$.
With energies scaled to $\sh^2/(2m\Delta r^2)$,
the interior tridiagonal matrix has
\begin{equation}
\boxed{H_{jj}=2+v_j,\quad
H_{j,j\pm1}=-1,\quad
v_j=-\frac{2\beta}{j}
+\frac{\ell(\ell+1)}{j^2}}.
\end{equation}
The end conditions are $u_0=u_{N+1}=0$.
For $\beta=1/50$, $N=1000$,
$b=(N+1)\aB/50\simeq20.02\aB$.
Diagonalizing the symmetric tridiagonal matrix,
and converting to Hartrees
$(1\ {\rm Ha}=\sh^2/(m\aB^2))$, gives
\begin{equation}
\begin{array}{c|rrr}
\ell& E_1/{\rm Ha}&E_2/{\rm Ha}&E_3/{\rm Ha}\\\hline
0&-0.499950&-0.124984&-0.049955\\
1&-0.124996&-0.051639&+0.007931\\
2&-0.053979&-0.005638&+0.067981
\end{array}.
\end{equation}
The corresponding unconfined hydrogen energies
are $E_n=-1/(2n^2)$ Ha, with
$n\ge\ell+1$. Higher-energy states are
strongly affected by the finite wall; in
particular positive confined eigenvalues
are not hydrogen bound levels.
The attached numerical script reproduces
the matrix eigenvalues and radial plots.
\note{The geometry here is a spherical \emph{volume}
with a hard radial boundary, not a surface-constrained
particle or a toroidal hydrogen system.}

\problem{4.69}{Numerical shooting targets for the first Coulomb states}
For the hydrogen radial equation, set
$E=-\sh^2/(2m\aB^2 n_{\rm eff}^2)$
and $\rho=2r/(n_{\rm eff}\aB)$.
The radial ODE is
\begin{equation}
\frac{\dd^2u}{\dd\rho^2}
+\left[-\frac14+\frac{n_{\rm eff}}{\rho}
-\frac{\ell(\ell+1)}{\rho^2}\right]u=0.
\end{equation}
Near $\rho=0$ the regular solution has
$u\propto\rho^{\ell+1}$.
Integrate from a small positive cutoff
$\rho_{\min}$ with
$u(\rho_{\min})=\rho_{\min}^{\ell+1}$
and $u'(\rho_{\min})=
(\ell+1)\rho_{\min}^\ell$.
For varying $n_{\rm eff}$, demand a
decaying rather than growing tail as
$\rho\to\infty$. The first three regular
normalizable values are
\begin{equation}
\boxed{\begin{array}{c|ccc}
\ell&n_{\rm eff,1}&n_{\rm eff,2}&n_{\rm eff,3}\\\hline
0&1.000&2.000&3.000\\
1&2.000&3.000&4.000\\
2&3.000&4.000&5.000
\end{array}}
\end{equation}
to four significant figures.
These values follow analytically from
termination of the confluent-hypergeometric
series at $n\ge\ell+1$. Shooting is
a numerical independent check; convergence
must be tested with larger integration
cutoffs and smaller step sizes.

\problem{4.70}{Sequential Stern--Gerlach measurements}
(a) For an initial $\ket{z+}$,
a first $z$ measurement always yields $z+$.
If the second measurement retains only
$x+$, every retained particle had $z+$
earlier. If instead the first measurement
was along $x$, the $x+$ and $x-$ outcomes
each occur with probability $1/2$,
but the second $x$ measurement retains only
particles that were previously $x+$.
Thus in either version
$\boxed{P(\text{first up}\mid\text{second }x+)=1}$.
(b) For arbitrary unit axes $\mathbf a,\mathbf b,
\mathbf c$, let
$u=\cos^2(\theta_{ab}/2)$,
$v=\cos^2(\theta_{bc}/2)$.
The joint probabilities for the two successive
measurements along $\mathbf b$ and $\mathbf c$
are
\begin{align}
P_{++}&=uv,&P_{+-}&=u(1-v),\\
P_{-+}&=(1-u)(1-v),&
P_{--}&=(1-u)v.
\end{align}
Among the subensemble whose second result is $c+$,
the fraction with an earlier $b+$ is
\begin{equation}
\boxed{\frac{uv}{uv+(1-u)(1-v)}
=\frac1{1+\tan^2(\theta_{ab}/2)
\tan^2(\theta_{bc}/2)}}.
\end{equation}

\problem{4.71}{Cartesian hydrogen $2p$ orbitals}
The combinations of spherical harmonics satisfy
\begin{align}
\sin\theta\cos\phi\,\sqrt{\frac{3}{4\pi}}
&=(Y_1^{-1}-Y_1^1)/\sqrt2,\\
\sin\theta\sin\phi\,\sqrt{\frac{3}{4\pi}}
&=\ii(Y_1^{-1}+Y_1^1)/\sqrt2,\\
\cos\theta\,\sqrt{\frac{3}{4\pi}}&=Y_1^0.
\end{align}
Multiplying by the common normalized $R_{21}(r)$,
\begin{equation}
\boxed{\psi_{2p_x}=
(\psi_{21,-1}-\psi_{211})/\sqrt2,\quad
\psi_{2p_y}=
\ii(\psi_{21,-1}+\psi_{211})/\sqrt2,\quad
\psi_{2p_z}=\psi_{210}}.
\end{equation}
These are proportional to
$x\ee^{-r/(2\aB)}$, $y\ee^{-r/(2\aB)}$,
$z\ee^{-r/(2\aB)}$, respectively.
Since, for example,
$L_x=-\ii\sh(y\partial_z-z\partial_y)$
and it annihilates any function
$x f(r)$, one has
\begin{equation}
\boxed{L_x\psi_{2p_x}=0,\quad
L_y\psi_{2p_y}=0,\quad
L_z\psi_{2p_z}=0}.
\end{equation}
Each aligned orbital is an eigenstate of
the corresponding angular-momentum component
with eigenvalue zero. Isodensity contours have
two lobes along the appropriate Cartesian axis,
separated by its perpendicular nodal plane.

\problem{4.72}{Uniform-field vector potential and magnetic coupling}
Let $\mathbf A=-\frac12\mathbf r\times\mathbf B_0
=\frac12\mathbf B_0\times\mathbf r$,
with constant $\mathbf B_0$.
The vector identity
$\nabla\times(\mathbf B_0\times\mathbf r)=
2\mathbf B_0$ gives
\begin{equation}
\boxed{\nabla\times\mathbf A=\mathbf B_0}.
\end{equation}
For a charge $q$ with spin coupling
$-\gamma\mathbf B_0\cdot\mathbf S$,
\begin{equation}
H=\frac{(\mathbf p-q\mathbf A)^2}{2m}
+q\phi-\gamma\mathbf B_0\cdot\mathbf S.
\end{equation}
Since $\nabla\cdot\mathbf A=0$,
$\mathbf p\cdot\mathbf A=\mathbf A\cdot\mathbf p$.
Also
$\mathbf A\cdot\mathbf p=
\frac12\mathbf B_0\cdot(\mathbf r\times\mathbf p)
=\frac12\mathbf B_0\cdot\mathbf L$
and
$\mathbf A^2=\frac14
[r^2B_0^2-(\mathbf r\cdot\mathbf B_0)^2]$.
Thus
\begin{equation}
\boxed{
H=\frac{p^2}{2m}+q\phi
-\mathbf B_0\cdot\left(\frac q{2m}\mathbf L+
\gamma\mathbf S\right)
+\frac{q^2}{8m}
[r^2B_0^2-(\mathbf r\cdot\mathbf B_0)^2]}.
\end{equation}
The term linear in field gives
the orbital/spin Zeeman interaction;
the quadratic term describes
diamagnetic coupling.

\problem{4.73}{Quantum force in a Stern--Gerlach field}
For a neutral spin-one-half particle
$H=p^2/(2m)-\gamma\mathbf B(\mathbf r)\cdot\mathbf S$.
With
$\mathbf B=-\alpha x\hat x+
(B_0+\alpha z)\hat z$,
only the spin interaction fails
to commute with $p_z$:
\begin{align}
\frac{\dd\avg{p_z}}{\dd t}
&=\frac{\ii}{\sh}
\avg{[H,p_z]}
=\gamma\avg{\partial_z\mathbf B\cdot\mathbf S}\\
&=\gamma\alpha\avg{S_z}.
\end{align}
Since $\dd\avg z/\dd t=\avg{p_z}/m$,
\begin{equation}
\boxed{m\frac{\dd^2\avg z}{\dd t^2}
=\gamma\alpha\avg{S_z}}.
\end{equation}
This is the quantum expectation-value
version of the Stern--Gerlach force law.

\problem{4.74}{Coupled spinor Schrödinger equations for Stern--Gerlach}
Write the complete state as
$\Psi=\Psi_+\chi_++\Psi_-\chi_-$,
and use
$\mathbf B=(-\alpha x,0,B_0+\alpha z)$.
The Hamiltonian matrix is
\begin{equation}
H=-\frac{\sh^2}{2m}\nabla^2I
-\frac{\gamma\sh}{2}
\begin{pmatrix}
B_0+\alpha z&-\alpha x\\
-\alpha x&-B_0-\alpha z
\end{pmatrix}.
\end{equation}
Equating spin components gives
\begin{align}
\ii\sh\partial_t\Psi_+
&=-\frac{\sh^2}{2m}\nabla^2\Psi_+
-\frac{\gamma\sh}{2}(B_0+\alpha z)\Psi_+
+\frac{\gamma\sh\alpha x}{2}\Psi_-,\\
\ii\sh\partial_t\Psi_-
&=-\frac{\sh^2}{2m}\nabla^2\Psi_-
+\frac{\gamma\sh}{2}(B_0+\alpha z)\Psi_-
+\frac{\gamma\sh\alpha x}{2}\Psi_+.
\end{align}
Factor the rapid uniform-field precession:
$\Psi_\pm=\ee^{\pm\ii\gamma B_0t/2}
\widetilde\Psi_\pm$. The equations become
\begin{align}
\ii\sh\partial_t\widetilde\Psi_+
&=\left[-\frac{\sh^2}{2m}\nabla^2
-\frac{\gamma\sh\alpha z}{2}\right]\widetilde\Psi_+
+\frac{\gamma\sh\alpha x}{2}
\ee^{-\ii\gamma B_0t}\widetilde\Psi_-,\\
\ii\sh\partial_t\widetilde\Psi_-
&=\left[-\frac{\sh^2}{2m}\nabla^2
+\frac{\gamma\sh\alpha z}{2}\right]\widetilde\Psi_-
+\frac{\gamma\sh\alpha x}{2}
\ee^{+\ii\gamma B_0t}\widetilde\Psi_+.
\end{align}
When rapid off-diagonal oscillations are
negligible, the spin components independently
see the linear potentials
$V_\pm=\mp\gamma\sh\alpha z/2$.
Their mean accelerations are
\begin{equation}
\boxed{\ddot{\avg z}_\pm
=\pm\frac{\gamma\sh\alpha}{2m}}.
\end{equation}
The packet splits into two paths.
\note{Neglecting off-diagonal coupling is
a controlled approximation only when the
uniform-field precession is rapid compared
with the spatial coupling dynamics.}

\problem{4.75}{Time-dependent magnetic flux in the ring problem}
For charge $q$ on a ring of radius $b$,
the time-dependent Hamiltonian is
\begin{equation}
H(t)=\frac{\sh^2}{2mb^2}
\left(-\ii\frac{\dd}{\dd\phi}
-\frac{q\Phi(t)}{2\pi\sh}\right)^2.
\end{equation}
All $H(t)$ are functions of the
same derivative operator and commute
at different times. A Fourier angular
mode $\ee^{\ii n\phi}/\sqrt{2\pi}$
has instantaneous eigenvalue
\begin{equation}
E_n(t)=\frac{\sh^2}{2mb^2}
\left(n-\frac{q\Phi(t)}{2\pi\sh}\right)^2.
\end{equation}
Let $f(t)=\sh^{-1}\int_0^tE_n(t')\dd t'$.
Then
\begin{equation}
\boxed{\Psi_n(\phi,t)=
\frac{\ee^{\ii n\phi}}{\sqrt{2\pi}}
\ee^{-\ii f(t)}}
\end{equation}
obeys
$\ii\sh\partial_t\Psi_n=
\sh f'(t)\Psi_n=E_n(t)\Psi_n=H(t)\Psi_n$.
Normalization is unchanged, as only
a time-dependent global phase is acquired.

\problem{4.76}{Mechanical angular momentum and changing flux}
Faraday's law gives the induced azimuthal
electric field on the ring:
$2\pi bE_\phi=-\dd\Phi/\dd t$.
For tangential speed $v=b\omega$,
the electric work rate is
\begin{equation}
\frac{\dd W}{\dd t}=qvE_\phi
=-\frac{q\omega}{2\pi}\frac{\dd\Phi}{\dd t},
\quad\boxed{\frac{\dd W}{\dd\Phi}
=-\frac{q\omega}{2\pi}}.
\end{equation}
In a state $\psi_n=\ee^{\ii n\phi}/\sqrt{2\pi}$,
canonical angular momentum is
$L_{z,\mathrm{can}}=n\sh$.
With $A_\phi=\Phi/(2\pi b)$, the
mechanical momentum is
$p_{\phi,\mathrm{mech}}=n\sh/b-qA_\phi$,
so
\begin{equation}
\boxed{L_{z,\mathrm{mech}}=
n\sh-\frac{q\Phi}{2\pi}},
\quad
\boxed{\omega=\frac{n\sh-q\Phi/(2\pi)}{mb^2}}.
\end{equation}
Differentiating the quantum stationary
energy at fixed integer $n$,
\begin{align}
\frac{\dd E_n}{\dd\Phi}
&=\frac{\dd}{\dd\Phi}
\frac{[n\sh-q\Phi/(2\pi)]^2}{2mb^2}\\
&=-\frac{q}{2\pi mb^2}
[n\sh-q\Phi/(2\pi)]\\
&=\boxed{-\frac{q\omega}{2\pi}}.
\end{align}
This exactly agrees with the
classical work rate per unit flux.
The \emph{mechanical} angular momentum
is not generally an integer multiple
of $\sh$ even though canonical
angular momentum is.
